% TEMPLATE-MILC-v3.tex - plantilla maestra UNICA de las guias MILC v3.
%
% La usan las tres familias de guias, cada una con su driver:
%   · Grados 6-11  -> scripts/build-guias-g11.py
%   · Bebras       -> scripts/build-guias-bebras.py
%   · Semillero    -> scripts/build-guias-semillero.py
%
% Antes habia tres copias casi identicas (~400 lineas iguales, 6 distintas):
% cada mejora habia que aplicarla tres veces, y en la practica no se hacia
% (el motor de recursos vivia solo en dos de ellas). Lo que varia entre
% familias esta parametrizado en 6 placeholders que cada driver llena
% (se nombran aqui SIN su sintaxis real: el builder reemplaza por texto plano,
% tambien dentro de los comentarios, y un valor multilinea rompe el preambulo):
%
%   HEADER_IZQ        encabezado izquierdo de pagina
%   HEADER_DER        encabezado derecho de pagina
%   PORTADA_ETIQUETA  rotulo sobre el numero grande (GRADO/MOMENTO/MODULO)
%   PORTADA_SERIE     linea de serie de la portada
%   PORTADA_ID        identificador de la guia en la portada
%   PORTADA_CIRCULO   el circulito de grado (°), o vacio si no aplica
%
% El resto de claves las documenta content/guias/_SCHEMA.md. El script valida
% que no queden placeholders sin reemplazar antes de escribir el archivo final.

\documentclass[11pt,letterpaper]{article}
\usepackage[margin=1.75cm]{geometry}
\usepackage[table]{xcolor}
\usepackage{fontspec}
\usepackage[spanish]{babel}
\usepackage{tikz}
% Librerías TikZ para los diagramas de las guías (bloque `recursos:`):
%  · babel  → babel en español vuelve activos caracteres como «>», y TikZ los usa
%             en sus opciones (>=stealth, ->). Sin esta librería, cualquier
%             diagrama con flechas revienta con «\language@active@arg> has an
%             extra }». Debe ir ANTES de que se dibuje nada.
%  · positioning        → sintaxis «below left=2pt and 3pt».
%  · decorations.pathreplacing → llaves ({brace}) para señalar rangos.
\usetikzlibrary{babel,positioning,decorations.pathreplacing}
\usepackage{graphicx}
% Circuitos: el trabajo de electrónica se hace hoy con micro:bit y sus sensores
% integrados (MakeCode, sin cableado), así que no se necesitan esquemáticos.
% Si algún día entran montajes con protoboard, basta descomentar esta línea:
% los diagramas .tex/.tikz declarados en `recursos:` ya se incrustan solos.
% \usepackage{circuitikz}
\usepackage[most]{tcolorbox}
\usepackage{tabularx}
\usepackage{array}
\usepackage{fancyhdr}
\usepackage{titlesec}
\usepackage{ragged2e}
\usepackage{enumitem}
\usepackage{microtype}

% Helvetica Neue no trae la flecha U+2192, asi que cada «→» escrito literalmente
% en el YAML se compilaba a NADA, en silencio (462 flechas en 61 guias: rutas del
% tipo «paleta → tipografia → lockup» salian rotas). Mapeamos el caracter a su
% version matematica, que si tiene glifo. Asi el autor puede seguir escribiendo
% «→» comodamente en el YAML.
\usepackage{newunicodechar}
\newunicodechar{→}{\ensuremath{\rightarrow}}

\setmainfont{Helvetica Neue}
\setsansfont{Helvetica Neue}
\setlength{\parindent}{0pt}
\setlength{\parskip}{6pt}
\hyphenpenalty=200
\exhyphenpenalty=200
\pretolerance=400
\tolerance=2000
\setlength{\emergencystretch}{1em}
\renewcommand{\arraystretch}{1.25}
\setlist{leftmargin=5mm,itemsep=1mm,topsep=1mm}
\newcolumntype{Y}{>{\RaggedRight\arraybackslash}X}

\definecolor{milcMagenta}{HTML}{E6007E}
\definecolor{milcVino}{HTML}{5A0038}
\definecolor{milcTurquesa}{HTML}{0093A5}
\definecolor{milcVerde}{HTML}{58A923}
\definecolor{milcMostaza}{HTML}{E5B400}
\definecolor{milcOcre}{HTML}{B86600}
\definecolor{milcNegro}{HTML}{241816}
\definecolor{milcGris}{HTML}{F7F7F4}
\definecolor{milcLinea}{HTML}{E8E2DD}
\definecolor{milcGradePrimary}{HTML}{FF2D87}
\definecolor{milcGradeDark}{HTML}{9F1B56}
\definecolor{milcGradeSoft}{HTML}{FFE0EE}
\definecolor{milcGradeAccent}{HTML}{CFE7FF}
\definecolor{milcGradeText}{HTML}{FFFFFF}
\color{milcNegro}

\pagestyle{fancy}
\fancyhf{}
\lhead{\small Tecnología e Informática · Grado 10}
\rhead{\small Guía 30 · MILC v3}
\cfoot{\small\thepage}
\renewcommand{\headrulewidth}{0pt}

\newtcolorbox{titlebox}[1]{
  enhanced,colback=#1,colframe=#1,arc=7mm,boxrule=0pt,
  left=7mm,right=7mm,top=3mm,bottom=3mm,
  before skip=8pt,after skip=8pt,
  fontupper=\sffamily\bfseries\Large\color{white}
}

\newtcolorbox{softbox}[3]{
  enhanced,colback=#2,colframe=#1,borderline west={4pt}{0pt}{#1},
  arc=2mm,boxrule=.4pt,left=4mm,right=4mm,top=3mm,bottom=3mm,
  before skip=6pt,after skip=8pt,title={#3},coltitle=white,
  colbacktitle=#1,fonttitle=\sffamily\bfseries,fontupper=\RaggedRight,
  attach boxed title to top left={xshift=3mm,yshift=-2mm},
  boxed title style={arc=2mm, boxrule=0pt}
}

\newtcolorbox{drawbox}[2]{
  enhanced,colback=white,colframe=#1,arc=4mm,boxrule=1pt,
  left=5mm,right=5mm,top=4mm,bottom=4mm,before skip=8pt,after skip=8pt,
  title={#2},coltitle=white,colbacktitle=#1,fonttitle=\sffamily\bfseries,
  fontupper=\RaggedRight,
  attach boxed title to top left={xshift=4mm,yshift=-2mm},
  boxed title style={arc=3mm, boxrule=0pt}
}

\newtcolorbox{infoband}[1]{
  enhanced,colback=#1!8,colframe=#1!50,arc=2mm,boxrule=.5pt,
  left=4mm,right=4mm,top=2mm,bottom=2mm,before skip=4pt,after skip=4pt,
  fontupper=\small\RaggedRight
}

% Puente narrativo entre secciones (contrato editorial Conéctate).
% Da continuidad: "Acabas de X. Ahora vas a Y porque Z."
% Visualmente: banda izquierda en color del grado, texto en itálica pequeña.
\newtcolorbox{puentebox}{
  enhanced,colback=milcGradeSoft,colframe=milcGradeDark,
  borderline west={3pt}{0pt}{milcGradeDark},
  arc=0mm,boxrule=0pt,left=5mm,right=4mm,top=2.5mm,bottom=2.5mm,
  before skip=4pt,after skip=8pt,
  fontupper=\itshape\small\color{milcNegro}\RaggedRight
}
% La flecha va en modo matematico a proposito: Helvetica Neue NO tiene el glifo
% U+2192, asi que \textrightarrow se compilaba a NADA («Missing character: There
% is no → in font Helvetica Neue») y el puente salia sin flecha. En modo
% matematico la flecha viene de la fuente matematica y si se dibuja.
\newcommand{\puente}[1]{%
  \begin{puentebox}%
    \textcolor{milcGradeDark}{$\rightarrow$}\hspace{2mm}#1%
  \end{puentebox}%
}

\newcommand{\linead}[1]{\noindent\color{milcLinea}\rule{#1}{0.5pt}\color{milcNegro}}
\newcommand{\respuesta}[1]{\par\vspace{2mm}\foreach \n in {1,...,#1}{\noindent\color{milcLinea}\rule{\linewidth}{0.5pt}\par\vspace{5mm}}\color{milcNegro}}
\newcommand{\checkbox}{\textcolor{milcGradeDark}{\fbox{\phantom{x}}}\hspace{5pt}}

\newcommand{\phasecard}[4]{
\begin{tcolorbox}[enhanced,colback=#3,colframe=#2,arc=3mm,boxrule=.5pt,height=3.05cm,valign=center,halign=center,left=1.5mm,right=1.5mm,top=2mm,bottom=2mm]
{\sffamily\bfseries\fontsize{22}{24}\selectfont\textcolor{#2}{#1}}\par
{\sffamily\bfseries\small #4}
\end{tcolorbox}
}

% ── Recursos visuales de la guía ──────────────────────────────────────
% Declarados en el bloque `recursos:` del YAML y emitidos por el builder.
% \guiaFigura   → imagen rasterizada o PDF (foto, carta celeste, montaje).
% \guiaDiagrama → fuente TikZ/circuitikz (.tex) incrustada como vector:
%                 es la vía para circuitos y esquemáticos editables.
% #1 = ruta relativa al .tex · #2 = pie (puede ir vacío)
\newcommand{\guiaFigura}[2]{%
\begin{center}
\includegraphics[width=0.88\linewidth,height=0.38\textheight,keepaspectratio]{#1}\\[1.5mm]
{\footnotesize\itshape\textcolor{milcNegro!75}{#2}}
\end{center}
}

\newcommand{\guiaDiagrama}[2]{%
\begin{center}
\input{#1}\\[1.5mm]
{\footnotesize\itshape\textcolor{milcNegro!75}{#2}}
\end{center}
}

\begin{document}
\thispagestyle{empty}
\begin{tikzpicture}[remember picture,overlay]
  \fill[white] (current page.south west) rectangle (current page.north east);
  \path[fill=milcGradePrimary,rounded corners=16mm]
    ([xshift=.55cm,yshift=-.55cm]current page.north west) rectangle
    ([xshift=-.55cm,yshift=3.65cm]current page.south east);

  \node[anchor=north west,text=milcGradeText] at ([xshift=2.0cm,yshift=-1.85cm]current page.north west)
    {\sffamily\bfseries\fontsize{14}{16}\selectfont GRADO};
  \node[anchor=north west,text=milcGradeText,opacity=.82] at ([xshift=2.0cm,yshift=-2.75cm]current page.north west)
    {\sffamily\bfseries\fontsize{9}{11}\selectfont SERIE GUÍAS · TECNOLOGÍA E INFORMÁTICA};

  \begin{scope}[shift={([xshift=-3.05cm,yshift=-2.55cm]current page.north east)}]
    \fill[white,opacity=.80] (-.62,-.52) rectangle (.62,.52);
    \draw[milcGradeText,line width=1pt] (-.62,-.52) rectangle (.62,.52);
    \fill[milcGradeDark] (-.42,-.38) rectangle (-.22,.06);
    \fill[milcGradeAccent] (-.10,-.38) rectangle (.10,.24);
    \fill[milcGradePrimary!60!black] (.22,-.38) rectangle (.42,.38);
  \end{scope}

  \node[anchor=west,text=milcGradeText] at ([xshift=1.9cm,yshift=-12.85cm]current page.north west)
    {\sffamily\bfseries\fontsize{124}{124}\selectfont 10};
  % El circulito de grado (°) solo aplica a las guias de grado («11°»); en
  % Bebras y Semillero el numero grande es un momento/modulo, no un grado.
  \draw[milcGradeText,line width=1.7pt]
    ([xshift=4.52cm,yshift=-11.68cm]current page.north west) circle (.13cm);

  \node[anchor=west,text=milcGradeText,text width=8.7cm,align=left] at ([xshift=10.35cm,yshift=-11.6cm]current page.north west)
    {
     {\sffamily\bfseries\fontsize{19}{22}\selectfont Sustentación final ---\\graduación del grado 10\\y feria de microempresas}\\[4mm]
     {\sffamily\bfseries\fontsize{23}{26}\selectfont Guía 30-10-TIC}\\[2mm]
     {\sffamily\fontsize{13}{16}\selectfont Período 3 · Ofimática con IA: contabilidad, datos y emprendimiento}\\[5mm]
     {\sffamily\bfseries\fontsize{11}{13}\selectfont Institución Educativa Sor María Juliana}\\[1mm]
     {\sffamily\fontsize{10}{12}\selectfont Autor: Dr. Álvaro Cárdenas Orozco}
    };

  \path[fill=milcGradeSoft,rounded corners=7mm]
    ([xshift=1.35cm,yshift=.85cm]current page.south west) rectangle
    ([xshift=-1.35cm,yshift=4.25cm]current page.south east);
  \draw[milcGradeDark,line width=.65pt,rounded corners=7mm]
    ([xshift=1.35cm,yshift=.85cm]current page.south west) rectangle
    ([xshift=-1.35cm,yshift=4.25cm]current page.south east);
  \node[anchor=center,text=milcNegro,text width=17.0cm,align=center] at ([yshift=2.55cm]current page.south)
    {\sffamily\fontsize{10}{12}\selectfont Metodología MILC · Apertura ancestral · Pensamiento computacional · Triángulo Dussel-Estoicismo-Floridi\\
    \textcolor{milcGradeDark}{\bfseries Producto:} Sustentación pública de 10 minutos del proyecto integrador frente a la comunidad escolar (compañeros, docentes y padres si es posible) + demo de los 3 productos digitales generados durante el año (libro de 80 páginas del P1, mini-proyecto del P2, micro-emprendimiento del P3) + carta abierta firmada de cierre del grado 10° + participación activa en feria de microempresas con stand del proyecto + autoevaluación final con compromiso para el grado 11°};
\end{tikzpicture}
\mbox{}
\newpage

\section*{Apertura: Sustentación final + feria de microempresas escolares}

\begin{softbox}{milcGradeDark}{milcGradeSoft}{Saber ancestral}
En los talleres antiguos del oficio, en las comunidades indígenas del Pacífico, en los gremios tradicionales europeos heredados al continente, hubo durante siglos un acto ritual que marcaba la transición del aprendiz al oficial: \textbf{la graduación frente a los mayores del oficio}. El aprendiz que había trabajado durante todo un periodo presentaba sus piezas maestras: las que demostraban que el oficio había sido aprendido. Esa presentación ocurría en público, en un día señalado, ante \textbf{el maestro, los oficiales viejos y los otros aprendices del taller}. El aprendiz exponía sus piezas, contaba cómo las hizo, declaraba dónde aún tenía dudas. Los mayores hacían preguntas técnicas y a veces incómodas. Si la defensa era sólida, recibía el \textbf{título del oficio}: ya no era aprendiz, era \emph{oficial}. La sabiduría era inquebrantable: \textbf{el oficio se certifica ante quienes lo conocen, no en privado}. Ese ritual ancestral es exactamente lo que la sustentación final del grado 10° + feria de microempresas escolares cierra: la graduación moderna que abre el camino al proyecto emprendedor final del grado 11°.
\end{softbox}

\begin{softbox}{milcTurquesa}{milcGris}{Saber contemporáneo}
La \textbf{sustentación final del grado 10° + feria de microempresas escolares} es el acto formal donde el estudiante presenta los 3 productos del año (libro de 80 páginas del P1, mini-proyecto del P2, micro-emprendimiento del P3) ante la comunidad escolar. Tiene 4 propiedades irrenunciables: (1) \textbf{Tiempo controlado}: 10 minutos de sustentación + 5 minutos Q\&A + tiempo en feria. (2) \textbf{Demostración real}: los 3 productos disponibles (PDF, archivos, Canvas impreso). (3) \textbf{Honestidad sobre uso de IA}: declarar abiertamente qué modelos usaste durante el año, qué porcentaje aproximado de los textos generaste con IA. (4) \textbf{Carta abierta firmada}: declaración personal del cierre del grado con compromiso explícito para el grado 11°. La \textbf{feria de microempresas} extiende la sustentación: cada estudiante pone stand con su micro-emprendimiento (incluso si es hipotético), los visitantes (compañeros + docentes + padres) recorren, conversan, intercambian comentarios. La regla profesional: \textbf{el cierre del grado 10° no es ejercicio escolar; es puente serio al proyecto emprendedor del grado 11°}.
\end{softbox}

\begin{softbox}{milcMostaza}{milcGris}{Pregunta-puente}
\textit{¿Qué sabía el aprendiz al exponer las piezas maestras ante los mayores del oficio, que el estudiante novato olvida cuando confunde cierre con \emph{``entregar tarea final''}? ¿Y por qué la feria de microempresas es la pieza más importante del cierre, aunque sea informal?}
\end{softbox}

\begin{softbox}{milcVerde}{milcGris}{Saber y saber hacer}
Al terminar podrás: (1) \textbf{analizar} los 3 productos del año identificando lo más representativo de cada periodo; (2) \textbf{explicar} en sustentación pública el viaje completo del año con honestidad técnica; (3) \textbf{evaluar} con honestidad preguntas duras de la audiencia; (4) \textbf{crear} la carta abierta firmada del cierre del grado con compromiso para el grado 11°.
\end{softbox}



\small
\begin{tabularx}{\textwidth}{>{\bfseries}p{3.0cm}Y}
\rowcolor{milcGradePrimary!12}Autor & Dr. Álvaro Cárdenas Orozco\\
DBA & Aplicación de ofimática asistida por IA a contabilidad básica y emprendimiento (MEN, T\&I)\\
Referentes & Lean Startup · Phronesis financiera · Ética del dato empresarial\\
Producto final & Sustentación pública de 10 minutos del proyecto integrador frente a la comunidad escolar (compañeros, docentes y padres si es posible) + demo de los 3 productos digitales generados durante el año (libro de 80 páginas del P1, mini-proyecto del P2, micro-emprendimiento del P3) + carta abierta firmada de cierre del grado 10° + participación activa en feria de microempresas con stand del proyecto + autoevaluación final con compromiso para el grado 11°\\
\end{tabularx}
\normalsize

\puente{Antes de empezar, mira el viaje completo. Has recorrido 29 sesiones del grado 10° en 3 periodos. Hoy es el \textbf{día de la graduación}. Después de hoy, cierras el grado 10° y abres el grado 11° (proyecto emprendedor final del bachillerato técnico).}

\begin{titlebox}{milcGradeDark}
Ruta MILC: así vamos a aprender
\end{titlebox}
\begin{tabularx}{\textwidth}{>{\centering\arraybackslash}X>{\centering\arraybackslash}X>{\centering\arraybackslash}X>{\centering\arraybackslash}X}
\phasecard{1}{milcVerde}{milcGris}{Escucha\\[-1mm]\small Observo, escucho y pregunto.} &
\phasecard{2}{milcTurquesa}{milcGris}{Sistematización\\[-1mm]\small Pensamiento computacional.} &
\phasecard{3}{milcMagenta}{milcGris}{Praxis\\[-1mm]\small Construyo y aplico.} &
\phasecard{4}{milcVino}{milcGris}{Evaluación\\[-1mm]\small Reflexión liberadora.}
\end{tabularx}


\puente{Antes de teorizar, vas a hacer una \emph{síntesis del año entero}: revisa los 3 productos (libro del P1, mini-proyecto del P2, micro-emprendimiento del P3) e identifica 1 cosa que estás orgulloso de cada uno + 1 cosa que mejorarías. Esa síntesis es la materia prima de la sustentación.}

\begin{titlebox}{milcVerde}
Fase 1 · Escucha
\end{titlebox}

\begin{softbox}{milcVerde}{milcGris}{Escena para escuchar}
\textbf{Actividad 1 · ANALIZA --- Síntesis del año en 3 productos} (15 min · individual). Revisa los 3 productos finales del año: (1) Libro de 80 páginas del P1. (2) Mini-proyecto del P2 (informe técnico con datos reales). (3) Micro-emprendimiento del P3 (proyecto integrador). Para cada uno, anota: (a) 1 cosa que estás orgulloso de haber logrado (lo que más representa tu aprendizaje). (b) 1 cosa que mejorarías si tuvieras más tiempo (lo que aún necesita revisión). Esa síntesis es la base de la sustentación.
\end{softbox}

\checkbox Revisé los 3 productos del año.\\
\checkbox Anoté 1 orgullo y 1 mejora por cada uno.\\
\checkbox Tengo síntesis del aprendizaje del año.

\begin{infoband}{milcVerde}


\textbf{Lo que va al cuaderno (Actividad 1 · ANALIZA).} Título: «Actividad 1 --- Síntesis del año». \textbf{Formato:} 3 fichas (1 por producto) con orgullo + mejora. \textbf{Sabes que terminaste cuando} podrías presentar los 3 con honestidad pública.
\end{infoband}

\puente{Ya tienes síntesis. Ahora vas a entender la \textbf{anatomía profesional} de la sustentación de 10 minutos + Q\&A + feria + carta abierta.}

\begin{titlebox}{milcTurquesa}
Fase 2 · Sistematización con pensamiento computacional
\end{titlebox}

\begin{softbox}{milcTurquesa}{milcGris}{Cuatro pilares aplicados al tema}
Tu síntesis es la base. Ahora necesitas la \textbf{anatomía profesional} de la sustentación + feria + carta.
\end{softbox}

\small
\begin{tabularx}{\textwidth}{>{\bfseries\RaggedRight\arraybackslash}p{3.6cm}Y}
\rowcolor{milcTurquesa!90}\color{white}Pilar & \color{white}\bfseries Aplicación al tema\\
1. Descomposición & La \textbf{estructura de la sustentación de 10 minutos} se descompone así (2 minutos por bloque): (1) \textbf{Introducción del año} (2 min): qué hicimos en el grado 10° (los 3 periodos), qué propósito tuvo. (2) \textbf{Libro del P1} (2 min): tema, proceso, fragmento corto leído en voz alta. (3) \textbf{Mini-proyecto del P2} (2 min): tema, datos clave, decisiones. (4) \textbf{Micro-emprendimiento del P3} (2 min): propuesta de valor, hallazgos del mercado, viabilidad. (5) \textbf{Declaración de uso de IA y carta abierta} (2 min): lectura de la carta firmada con compromiso para el grado 11°. Después: \textbf{Q\&A de 5 minutos} con preguntas de la audiencia + \textbf{feria de microempresas} con stand.\\
2. Reconocimiento de patrones & La \textbf{feria de microempresas} es la pieza más activa del cierre. Cada estudiante pone stand con su micro-emprendimiento del P3 (incluso si es hipotético): impresión de portada del libro, copia del Canvas, propuesta comercial. Los visitantes (compañeros + docentes + padres si pueden venir) recorren, conversan, hacen preguntas, dejan comentarios. Es práctica de comunicación profesional en contexto realista. La regla: \textbf{cada autor está 30 min en su stand y 30 min recorriendo otros}.\\
3. Abstracción & La \textbf{carta abierta firmada} es la pieza más personal del cierre. Tiene 5 párrafos: (a) \textbf{Apertura}: qué fue para mí este año en el grado 10°. (b) \textbf{Lo que aprendí del P1}: 1 aprendizaje concreto. (c) \textbf{Lo que aprendí del P2}: 1 aprendizaje concreto. (d) \textbf{Lo que aprendí del P3}: 1 aprendizaje concreto. (e) \textbf{Compromiso para el grado 11°}: 1 acción que me llevo. Firmada con nombre completo + fecha + ciudad. Leída en voz alta durante la sustentación.\\
4. Algoritmo conceptual & Algoritmo: (1) preparar 5 slides (1 por bloque). (2) escribir guion hablado de 1500 palabras (= 10 min). (3) ensayar cronometrado. (4) sustentar. (5) responder Q\&A con honestidad. (6) participar en feria como autor y como visitante. (7) leer carta abierta. (8) autoevaluar.\\
\end{tabularx}
\normalsize

\begin{drawbox}{milcTurquesa}{Anatomía del cierre del grado 10}
\textbf{Sustentación 10 min:} intro + 3 productos + carta. \\[1mm] \textbf{Q\&A 5 min:} honestidad ante preguntas. \\[1mm] \textbf{Feria:} stand + recorrido a otros. \\[1mm] \textbf{Carta abierta:} 5 párrafos firmados.
\end{drawbox}

\begin{softbox}{milcMostaza}{milcGris}{Errores comunes y su corrección}
(1) Leer las slides. (2) Ocultar uso de IA. (3) Carta genérica sin compromiso específico. (4) Sin stand preparado para la feria. (5) Visitar solo amigos cercanos en la feria. (6) Pasarse del tiempo. (7) Cierre apurado sin proyección al grado 11°.
\end{softbox}

\begin{infoband}{milcTurquesa}


\textbf{Lo que va al cuaderno (Actividad 2 · EXPLICA).} Título: «Actividad 2 --- Anatomía cierre». \textbf{Formato:} (a) estructura 10 min; (b) reglas de la feria; (c) plantilla carta abierta; (d) los 7 errores con marca. \textbf{Sabes que terminaste cuando} tienes plan claro de la sesión.
\end{infoband}

\puente{Tienes el método. Toca aplicarlo: preparas slides simples, ensayas, sustentas, participas en feria, lees carta abierta firmada, autoevaluación final.}

\begin{titlebox}{milcMagenta}
Fase 3 · Praxis con pensamiento computacional
\end{titlebox}

\begin{softbox}{milcMagenta}{milcGris}{Construyendo el producto con los cuatro pilares}
\textbf{Actividad 3 · EVALÚA + CREA --- Sustentación + feria + carta + autoevaluación} (90 min total · individual + colectivo). Hora de ejecutar. Preparas slides, ensayas, sustentas, participas en feria, lees carta firmada, autoevalúas con compromiso.
\end{softbox}

\small
\begin{tabularx}{\textwidth}{>{\bfseries\RaggedRight\arraybackslash}p{3.6cm}Y}
\rowcolor{milcMagenta!90}\color{white}Pilar & \color{white}\bfseries Aplicación al producto\\
Descomposición & Tu cierre se construye en 8 pasos: (1) preparar 5 slides simples. (2) escribir guion hablado de 1500 palabras. (3) ensayar cronometrado. (4) preparar stand para la feria (impresiones, Canvas, productos). (5) sustentar en público con honestidad técnica. (6) responder Q\&A. (7) participar en feria 30 min como autor y 30 min recorriendo. (8) leer carta abierta firmada y autoevaluar con compromiso para el grado 11°.\\
Patrones & Sigue el patrón \textbf{``orgullo honesto y mejora declarada''}: el cierre del año no es momento de inflar ni de subestimar. Reconoce con orgullo lo que lograste (\emph{``terminé el libro de 80 páginas''}, \emph{``apliqué encuesta a 18 personas''}) y declara con honestidad lo que aún te falta (\emph{``mi Canvas tiene supuestos sin verificar''}, \emph{``mi correo profesional aún suena formal''}). Ambos lados son virtud.\\
Abstracción & La \textbf{carta abierta firmada} es tu firma del año. Modelo de párrafo de cierre con compromiso: \emph{``Me llevo del grado 10° la habilidad de trabajar con IA respetando autoría. Me comprometo para el grado 11° a aplicar todo lo aprendido en proyecto emprendedor real con destinatario verificable, no hipotético''}. Específico, no genérico.\\
Algoritmo de armado & Algoritmo: (1) slides. (2) guion. (3) ensayo. (4) stand. (5) sustentación. (6) Q\&A. (7) feria. (8) carta abierta. (9) autoevaluación. (10) entrega al docente: link a los 3 productos + slides + guion + carta firmada.\\
\end{tabularx}
\normalsize

\begin{drawbox}{milcMagenta}{Checklist antes de sustentar}
\checkbox 5 slides simples preparadas. \\[1mm] \checkbox Guion hablado de 1500 palabras. \\[1mm] \checkbox Ensayo cronometrado hecho. \\[1mm] \checkbox Stand de feria preparado. \\[1mm] \checkbox Carta abierta de 5 párrafos firmada. \\[1mm] \checkbox Declaración de uso de IA del año preparada. \\[1mm] \checkbox Compromiso para el grado 11° específico.
\end{drawbox}

\begin{softbox}{milcMagenta}{milcGris}{Plantilla de trabajo}
\textbf{2 min Intro.} \textit{Año en 3 periodos.} \\[1mm] \textbf{2 min Libro P1.} \textit{Tema + fragmento corto.} \\[1mm] \textbf{2 min Informe P2.} \textit{Datos + decisiones.} \\[1mm] \textbf{2 min Micro-emprendimiento P3.} \textit{Propuesta + viabilidad.} \\[1mm] \textbf{2 min Carta abierta.} \textit{Lectura + firma.} \\[1mm] \textbf{Q\&A 5 min.} \textit{Honestidad.} \\[1mm] \textbf{Feria.} \textit{Stand + recorrido.}
\end{softbox}

\begin{infoband}{milcMagenta}


\textbf{Lo que va al cuaderno (Actividad 3 · EVALÚA + CREA).} Título: «Actividad 3 --- Mi cierre del grado 10». \textbf{Formato:} (a) slides; (b) guion; (c) carta abierta firmada; (d) autoevaluación + compromiso para G11. \textbf{Sabes que terminaste cuando} entregaste materiales y tomaste tiempo para reflexionar.
\end{infoband}

\puente{Lo que sigue resume \emph{qué} entregas hoy: sustentación ejecutada + feria participada + carta firmada + autoevaluación con compromiso. El criterio principal: que la audiencia se vaya sabiendo qué construiste durante el año y para qué.}

\begin{titlebox}{milcMostaza}
Producto de la sesión
\end{titlebox}
\textbf{Sustentación 10 min + 3 productos + feria + carta firmada + autoevaluación}.
\begin{softbox}{milcMostaza}{milcGris}{Criterios de calidad}
(1) Sustentación dentro de 10 minutos. (2) 5 slides simples por bloque. (3) 3 productos del año mostrados. (4) Q\&A con honestidad técnica. (5) Carta abierta de 5 párrafos firmada. (6) Participación activa en feria como autor y visitante. (7) Compromiso específico para grado 11.
\end{softbox}

\puente{Antes de cerrar, mira la sustentación desde las cinco dimensiones humanas. Cerrar el grado con honestidad y orgullo es respeto por las 29 sesiones; cerrar sin asumir lo que falta es traicionar al estudiante que fuiste durante el año.}

\begin{titlebox}{milcVino}
Fase 4 · Evaluación liberadora — cinco dimensiones
\end{titlebox}

\begin{softbox}{milcVino}{milcGris}{Sentido de la evaluación}
Evaluar no es solo revisar una nota. Es reconocer qué aprendí, qué cambió en mí, qué emociones manejé, cómo me relaciono con la ciudad y la comunidad, y qué le dejaré a quien viene detrás.
\end{softbox}

\textbf{Mi reflexión liberadora en cinco dimensiones.} Completa cada frase iniciada:

\small
\begin{tabularx}{\textwidth}{>{\bfseries\RaggedRight\arraybackslash}p{3.4cm}Y>{\RaggedRight\arraybackslash}p{4.6cm}}
\rowcolor{milcVino}\color{white}Dimensión & \color{white}\bfseries Mi respuesta (frame starter) & \color{white}\bfseries Algo que descubrí\\
1. Personal & Lo que más cambió en mí fue \linead{6.5cm} & \linead{4.2cm}\\
2. Emocional & La emoción más fuerte fue \linead{2.5cm} y la regulé \linead{3cm} & \linead{4.2cm}\\
3. Ciudadana & En democracia esto importa porque \linead{6cm} & \linead{4.2cm}\\
4. Local (Cartago) & En mi barrio sería útil para \linead{6.5cm} & \linead{4.2cm}\\
5. Intergeneracional & A mi abuela/abuelo le diría \linead{6.5cm} & \linead{4.2cm}\\
\end{tabularx}
\normalsize

\begin{infoband}{milcVino}
\textbf{Para la próxima sustentación.} Vuelve a esta tabla en seis meses. Verás que las respuestas cambian — no porque lo aprendido cambie, sino porque tú cambias. La evaluación liberadora es una conversación contigo a través del tiempo.
\end{infoband}


\puente{Y para terminar, tres voces sobre el cierre del oficio. Dussel sobre el aprendiz que devuelve a la comunidad lo aprendido, Marco Aurelio sobre la honestidad del cierre, Floridi sobre la graduación con declaración de IA como modelo profesional emergente.}

\begin{titlebox}{milcGradeDark}
Triángulo de pensamiento · cierre filosófico
\end{titlebox}

\begin{softbox}{milcVino}{milcGris}{Dussel · lente del nosotros}
\textit{``El aprendiz que devuelve oficio asumido a la comunidad cierra el ciclo de la pedagogía liberadora.''}

Tu sustentación cierra el grado 10°: lo que aprendiste vuelve a la comunidad escolar hecho voz propia firmada. La filosofía de la liberación pide ese acto. Cuando declaras con honestidad lo que lograste y lo que aún te falta, estás devolviendo oficio asumido. La phronesis del cierre consiste en \textbf{compartir el aprendizaje del año, no guardarlo}.

\textbf{¿Mi sustentación y mi participación en feria devuelven oficio a la comunidad escolar, o solo cumplen con requisito final?}
\end{softbox}
\linead{\linewidth}\\[1mm]
\linead{\linewidth}

\begin{softbox}{milcMostaza}{milcGris}{Estoicismo · Marco Aurelio · cuidado interior}
\textit{``Reconocer con honestidad lo aprendido y lo que falta es virtud final del año.''}

Es tentador presentar el año como perfecto cuando hubo dificultades. La disciplina estoica recomienda lo opuesto: \emph{declara con honestidad lo que aprendiste y lo que aún te cuesta}. Esa precisión gana respeto. La phronesis del cierre honra la verdad sobre la imagen.

\textbf{¿Mi sustentación declara honestamente lo que aprendí y lo que aún me falta?}
\end{softbox}
\linead{\linewidth}\\[1mm]
\linead{\linewidth}

\begin{softbox}{milcTurquesa}{milcGris}{Floridi · lente de la infoesfera}
\textit{``La graduación con declaración honesta de IA es modelo profesional emergente del siglo XXI.''}

En la era de la IA generativa, los cierres profesionales que declaran abiertamente uso de IA durante el año se posicionan como ético contemporáneo. Tu cierre del grado 10° te entrena en esa práctica que la educación profesional adoptará obligatoriamente en pocos años: \textbf{transparencia con porcentajes y compromiso de mejora}.

\textbf{¿Mi cierre contribuye al modelo de graduación transparente del siglo XXI?}
\end{softbox}
\linead{\linewidth}\\[1mm]
\linead{\linewidth}

\begin{titlebox}{milcVino}
Compromiso final
\end{titlebox}

\small
\begin{tabularx}{\textwidth}{>{\RaggedRight\arraybackslash}p{3.0cm}YYY}
\rowcolor{milcVino}\color{white}\bfseries Criterio & \color{white}\bfseries Lo logré & \color{white}\bfseries En proceso & \color{white}\bfseries Necesito apoyo\\
Comprensión del tema & Explico ideas centrales con mis palabras. & Reconozco ideas pero falta precisión. & Necesito repasar conceptos base.\\
Pensamiento computacional & Apliqué los 4 pilares al diseñar mi producto. & Apliqué algunos pilares. & Necesito releer la fase 2.\\
Aplicación y producto & Mi producto cumple los criterios y se puede revisar. & Producto funcional con ajustes pendientes. & Aún en construcción.\\
Reflexión 5 dimensiones & Respondí cada dimensión con honestidad. & Respondí algunas con detalle. & Mis respuestas son muy generales.\\
Triángulo de pensamiento & Conecté las 3 voces con mi trabajo real. & Conecté al menos una. & No relaciono las voces con mi proceso.\\
\end{tabularx}
\normalsize

\textbf{Hoy entendí que:}\\[1mm]
\linead{\linewidth}\\[1mm]
\linead{\linewidth}

\textbf{Mi compromiso para la próxima sesión es:}\\[1mm]
\linead{\linewidth}\\[1mm]
\linead{\linewidth}

\begin{softbox}{milcMostaza}{milcGris}{Cita de despedida · Marco Aurelio}
\textit{``El obstáculo en el camino se vuelve el camino. Lo que estorba al camino se vuelve camino.''}\\[1mm]
Cada error de hoy, bien observado, se convierte en pieza del camino que pisarás mañana.
\end{softbox}

\small
\begin{tabularx}{\textwidth}{>{\bfseries}p{3.6cm}Y}
\rowcolor{milcGradeDark!12}Nombre completo: & \linead{10cm}\\
Fecha: & \linead{4cm} \quad Curso/grupo: \linead{4cm}\\
Mi firma: & \linead{9cm}\\
\end{tabularx}
\normalsize

\begin{infoband}{milcGradeDark}
\textbf{Después del grado 10°.} Tu carta abierta firmada es el cierre formal del grado, pero también el primer documento que llevas al grado 11°. Guarda copia personal. Al entrar al grado 11° (proyecto emprendedor final del bachillerato técnico), relee la carta: tu compromiso para el grado 11° es lo primero que aplicarás. El grado 10° preparó las herramientas; el grado 11° es donde se aplican a un proyecto emprendedor real con destinatario verificable.
\end{infoband}

\end{document}
